\chapter*{ЗАКЛЮЧЕНИЕ}
\addcontentsline{toc}{chapter}{ЗАКЛЮЧЕНИЕ}

Цель работы достигнута. Разработана рекомендательная система ликёров, объединяющая контентные рекомендации и параметрический поиск. Все проектные решения обоснованы сравнением альтернатив.

В ЛР1 сформирован датасет из 30 ликёров с признаками пяти типов и иерархия предметной области по вкусовой основе. В ней 42 узла классификации при высоте 5 уровней. Деление по вкусу выбрано как смыслообразующее для покупателя. Числовые признаки оставлены атрибутами объектов.

В ЛР2 для каждого типа признаков выбрана частная мера по формальной методике. Она включает матрицы смежности, тепловые карты, эталонные пары, корреляцию ранжирований и зернистость шкалы. Выбраны евклидова мера, SMC, мера У--Палмера, порядковая мера и $\delta$-функция. Обобщающая мера --- взвешенная сумма с экспертными весами. Её устойчивость проверена пятью объективными схемами весов.

В ЛР3 разработана контентная РС в трёх версиях. Для набора лайков выбран подход <<группы лайков со сшиванием>> с $\theta = 0{,}65$. При трёх интересах recall(5) равен 0,797, а покрытие интересов --- 100\,\%. Для сравнения у среднего эти значения равны 0,649 и 74\,\%. Для дизлайков выбран контрастный штраф. Он втрое снижает утечку нелюбимого стиля и сохраняет любимый жанр. Каждая рекомендация объясняется.

В ЛР4 реализован поиск по десяти фильтрам с проверкой ввода. Пустая выдача заменяется ранжированием по функции полезности. Задуманный объект попадает в первую пятёрку во всех 179 сценариях, а при ослаблении фильтров --- только в 70\,\%. Ослабление используется как подсказка. В РК1 поиск и рекомендации объединены в одну систему с общим интерфейсом, сессионным профилем, фильтрацией рекомендаций и историей шагов.
